%【紙面設定】
%b4,横書き,2段
\documentclass[paper=b4j,landscape,twocolumn,fleqn]{jlreq}
%横書き用の調整
\special{papersize=\the\paperwidth,\the\paperheight}
%間に線をいれる
\setlength{\columnseprule}{0.4pt}
%ページ番号消す
\pagestyle{empty}
%ページ余白の調整
\usepackage[margin=8truemm]{geometry}
%画像挿入用のパッケージ
\usepackage[dvipdfmx]{graphicx}
%数学用のパッケージ
\usepackage{amsmath}
%見出しによる行取りの数を減らす
\ModifyHeading{section}{
	lines=1,
	before_space=6pt,
	after_space=0pt
}
% 段落・数式まわりの余白を詰める
\setlength{\parskip}{0pt}
\setlength{\jot}{3pt}

%【本文】
\begin{document}
\normalsize
%1段落目
\section{覚えるべき三角形の図示}
次の各問について、図を描き、辺の長さの比を記入せよ。
\begin{align*}
　&1.　30^\circ-60^\circ-90^\circ\text{ の直角三角形を図示し、辺の比を記入せよ。}\\[0.45em]
　&2.　45^\circ-45^\circ-90^\circ\text{ の直角三角形を図示し、辺の比を記入せよ。}\\[0.45em]
　&3.　30^\circ-60^\circ-90^\circ\text{ の三角形で、最短辺を1とした図を描き、他の2辺を記入せよ。}\\[0.45em]
　&4.　45^\circ-45^\circ-90^\circ\text{ の三角形で、斜辺を1とした図を描き、他の2辺を記入せよ。}\\[0.45em]
　&5.　30^\circ,45^\circ,60^\circ\text{ について、}sin\theta,\ cos\theta\text{ の値を三角形の図から読み取り記入せよ。}\\[0.45em]
　&6.　30^\circ,45^\circ,60^\circ\text{ について、}tan\theta\text{ の値を三角形の図から読み取り記入せよ。}\\[0.45em]
\end{align*}

\section{直角三角形の比}
直角三角形について、以下の問いに答えよ。
\begin{align*}
　&1.　30^\circ-60^\circ-90^\circ の直角三角形で、最も短い辺が2のとき、残り2辺の長さを求めよ。\\[0.45em]
　&2.　30^\circ-60^\circ-90^\circ の直角三角形で、斜辺が8のとき、残り2辺の長さを求めよ。\\[0.45em]
　&3.　45^\circ-45^\circ-90^\circ の直角三角形で、斜辺が6\sqrt{2}のとき、残り2辺の長さを求めよ。\\[0.45em]
　&4.　60^\circ に対する向かい側の辺が3\sqrt{3}のとき、30^\circ に対する向かい側の辺と斜辺を求めよ。\\[0.45em]
　&5.　30^\circ に対する向かい側の辺が5のとき、60^\circ に対する向かい側の辺と斜辺を求めよ。\\[0.45em]
　&6.　45^\circ-45^\circ-90^\circ の直角三角形で、1辺が4のとき、もう1辺と斜辺を求めよ。\\[0.45em]
\end{align*}

\section{三角比}
直角三角形において、角$\theta$に対する三角比を考える。
\begin{align*}
　&1.　\cos 30^\circ,\ \cos 45^\circ,\ \cos 60^\circ\text{ の値を求めよ。}\\[0.45em]
　&2.　\sin 30^\circ,\ \sin 45^\circ,\ \sin 60^\circ\text{ の値を求めよ。}\\[0.45em]
　&3.　\tan 30^\circ,\ \tan 45^\circ,\ \tan 60^\circ\text{ の値を求めよ。}\\[0.45em]
　&4.　斜辺が8、\theta=30^\circ のとき、底辺と高さを求めよ。\\[0.45em]
　&5.　斜辺が10、\theta=60^\circ のとき、底辺と高さを求めよ。\\[0.45em]
　&6.　底辺が6、\theta=45^\circ のとき、斜辺と高さを求めよ。\\[0.45em]
\end{align*}

\newpage

\section{覚えるべき三角形の図示}
次の各問について、図を描き、辺の長さの比を記入せよ。
\begin{align*}
　&1.　30^\circ-60^\circ-90^\circ\text{ の直角三角形を図示し、辺の比を記入せよ。}\\[0.45em]
　&2.　45^\circ-45^\circ-90^\circ\text{ の直角三角形を図示し、辺の比を記入せよ。}\\[0.45em]
　&3.　30^\circ-60^\circ-90^\circ\text{ の三角形で、最短辺を1とした図を描き、他の2辺を記入せよ。}\\[0.45em]
　&4.　45^\circ-45^\circ-90^\circ\text{ の三角形で、斜辺を1とした図を描き、他の2辺を記入せよ。}\\[0.45em]
　&5.　30^\circ,45^\circ,60^\circ\text{ について、}sin\theta,\ cos\theta\text{ の値を三角形の図から読み取り記入せよ。}\\[0.45em]
　&6.　30^\circ,45^\circ,60^\circ\text{ について、}tan\theta\text{ の値を三角形の図から読み取り記入せよ。}\\[0.45em]
\end{align*}

\section{直角三角形の比}
直角三角形について、以下の問いに答えよ。
\begin{align*}
　&1.　30^\circ-60^\circ-90^\circ の直角三角形で、最も短い辺が2のとき、残り2辺の長さを求めよ。\\[0.45em]
　&2.　30^\circ-60^\circ-90^\circ の直角三角形で、斜辺が8のとき、残り2辺の長さを求めよ。\\[0.45em]
　&3.　45^\circ-45^\circ-90^\circ の直角三角形で、斜辺が6\sqrt{2}のとき、残り2辺の長さを求めよ。\\[0.45em]
　&4.　60^\circ に対する向かい側の辺が3\sqrt{3}のとき、30^\circ に対する向かい側の辺と斜辺を求めよ。\\[0.45em]
　&5.　30^\circ に対する向かい側の辺が5のとき、60^\circ に対する向かい側の辺と斜辺を求めよ。\\[0.45em]
　&6.　45^\circ-45^\circ-90^\circ の直角三角形で、1辺が4のとき、もう1辺と斜辺を求めよ。\\[0.45em]
\end{align*}

\section{三角比}
直角三角形において、角$\theta$に対する三角比を考える。
\begin{align*}
　&1.　\cos 30^\circ,\ \cos 45^\circ,\ \cos 60^\circ\text{ の値を求めよ。}\\[0.45em]
　&2.　\sin 30^\circ,\ \sin 45^\circ,\ \sin 60^\circ\text{ の値を求めよ。}\\[0.45em]
　&3.　\tan 30^\circ,\ \tan 45^\circ,\ \tan 60^\circ\text{ の値を求めよ。}\\[0.45em]
　&4.　斜辺が8、\theta=30^\circ のとき、底辺と高さを求めよ。\\[0.45em]
　&5.　斜辺が10、\theta=60^\circ のとき、底辺と高さを求めよ。\\[0.45em]
　&6.　底辺が6、\theta=45^\circ のとき、斜辺と高さを求めよ。\\[0.45em]
\end{align*}
\end{document}



